\chapter{Permutations and Combinations}
\signature

\section{The basic counting rules}

\begin{theorem}[Sum rule]
If a task can be done in $n_1$ ways or, \emph{exclusively}, in $n_2$ ways,
it can be done in $n_1 + n_2$ ways. Set version: disjoint sets satisfy
$|A \cup B| = |A| + |B|$.
\end{theorem}

\begin{theorem}[Product rule]
If a procedure consists of two successive independent choices, with $n_1$
options for the first and $n_2$ for the second, there are $n_1 n_2$ outcomes.
Set version: $|A \times B| = |A|\,|B|$.
\end{theorem}

\begin{example}
A password of $3$ letters (26 choices each) followed by $2$ digits:
$26^3 \cdot 10^2 = 1{,}757{,}600$. A menu offering $4$ starters or $6$ mains
as a single dish: $4 + 6 = 10$ choices. Tree diagrams display such combined
choices as branches; the leaves count the outcomes.
\end{example}

\section{Permutations}

\begin{definition}
A \emph{permutation} of $n$ distinct objects is an ordering of all of them.
The number of permutations is
$n! = n(n-1)(n-2)\cdots 2 \cdot 1$, with $0! = 1$.
\end{definition}

\begin{theorem}[$k$-permutations]
The number of ordered arrangements of $k$ objects chosen from $n$ distinct
objects is
\[
P(n,k) = n(n-1)\cdots(n-k+1) = \frac{n!}{(n-k)!}.
\]
\end{theorem}

\begin{example}[Sprint podium]
Nine athletes run a sprint; podium orders (gold, silver, bronze) number
$P(9,3) = 9 \cdot 8 \cdot 7 = 504$.
\end{example}

\begin{theorem}[Permutations with repeated objects]\label{thm:msperm}
The number of distinct arrangements of $n$ objects of which $n_1$ are alike
of type 1, \ldots, $n_k$ alike of type $k$ ($n_1 + \cdots + n_k = n$) is
\[
\frac{n!}{n_1!\,n_2!\cdots n_k!}.
\]
\end{theorem}

\begin{example}
``MISSISSIPPI'' has $\frac{11!}{1!\,4!\,4!\,2!} = 34650$ distinct
arrangements.
\end{example}

\begin{theorem}[Circular permutations]
$n$ distinct people can be seated around a round table (rotations considered
identical) in $(n-1)!$ ways.
\end{theorem}

\section{Combinations}

\begin{definition}
A \emph{combination} of $k$ objects from $n$ is an unordered selection. Their
number is the \emph{binomial coefficient}
\[
C(n,k) = \binom{n}{k} = \frac{n!}{k!\,(n-k)!} = \frac{P(n,k)}{k!}.
\]
\end{definition}

\begin{example}
Choosing a committee of $5$ from $30$ students:
$\binom{30}{5} = 142506$. Note order is irrelevant --- choosing
$\{$Amal, Badr$\}$ equals choosing $\{$Badr, Amal$\}$.
\end{example}

\begin{theorem}[Symmetry and Pascal's identity]\label{thm:pascal}
\[
\binom{n}{k} = \binom{n}{n-k},
\qquad
\binom{n+1}{k} = \binom{n}{k-1} + \binom{n}{k}.
\]
\end{theorem}

\begin{proof}
Symmetry: choosing the $k$ selected is the same as choosing the $n-k$
rejected. Pascal: fix one element $x$ of an $(n+1)$-set; the $k$-subsets
split into those containing $x$ ($\binom{n}{k-1}$ of them) and those not
($\binom{n}{k}$).
\end{proof}

Pascal's identity generates \emph{Pascal's triangle}, whose row $n$ lists
$\binom{n}{0}, \ldots, \binom{n}{n}$:
\[
\begin{array}{ccccccccc}
&&&&1&&&&\\
&&&1&&1&&&\\
&&1&&2&&1&&\\
&1&&3&&3&&1&\\
1&&4&&6&&4&&1
\end{array}
\]
Row sums give $\sum_k \binom{n}{k} = 2^n$ (each subset of an $n$-set has some
size $k$ --- recounting the power-set theorem of Chapter~6).

\section{The binomial theorem}

\begin{theorem}[Binomial theorem]\label{thm:binom}
For all $x, y$ and $n \in \N$:
\[
(x + y)^{n} = \sum_{k=0}^{n} \binom{n}{k} x^{\,n-k} y^{k}.
\]
\end{theorem}

\begin{proof}
Expanding $(x+y)(x+y)\cdots(x+y)$, each term arises by choosing $y$ from $k$
of the $n$ factors and $x$ from the rest; there are $\binom{n}{k}$ such
choices, giving the coefficient of $x^{n-k}y^{k}$.
\end{proof}

\section{Combinations with repetition}

\begin{theorem}[Stars and bars]\label{thm:stars}
The number of multisets of size $k$ from $n$ types --- equivalently the
number of solutions of $x_1 + \cdots + x_n = k$ in nonnegative integers ---
is
\[
\binom{n + k - 1}{k}.
\]
\end{theorem}

\begin{proof}
Encode a selection as $k$ stars separated by $n-1$ bars; a word with
$n + k - 1$ symbols is determined by the positions of the $k$ stars.
\end{proof}

\subsection{The four counting problems}

\begin{center}
\renewcommand{\arraystretch}{1.4}
\begin{tabular}{lcc}
\toprule
Select $k$ from $n$ & order matters & order irrelevant\\
\midrule
without repetition & $P(n,k) = \dfrac{n!}{(n-k)!}$ & $\dbinom{n}{k}$\\
with repetition & $n^{k}$ & $\dbinom{n+k-1}{k}$\\
\bottomrule
\end{tabular}
\end{center}

\section{Multinomial coefficients}

\begin{theorem}[Multinomial theorem]
\[
(x_1 + \cdots + x_m)^{n}
= \sum_{n_1 + \cdots + n_m = n} \binom{n}{n_1, n_2, \ldots, n_m}\,
x_1^{n_1} \cdots x_m^{n_m},
\qquad
\binom{n}{n_1, \ldots, n_m} = \frac{n!}{n_1! \cdots n_m!},
\]
the multinomial coefficient counting the ways to partition $n$ labelled
objects into groups of sizes $n_1, \ldots, n_m$ (cf.\
Theorem~\ref{thm:msperm}).
\end{theorem}

\section{The pigeonhole principle}

\begin{theorem}[Pigeonhole principle]\label{thm:pigeon}
If $n+1$ or more objects are placed into $n$ boxes, some box contains at
least two objects. Generalised form: with $N$ objects in $n$ boxes, some box
contains at least $\lceil N/n \rceil$ objects.
\end{theorem}

\begin{example}
Among any $13$ people, two share a birth month. Among $100$ people, at least
$\lceil 100/12 \rceil = 9$ share one. Subtler: among any $n+1$ integers
chosen from $\{1, \ldots, 2n\}$, two are consecutive --- pair the numbers
$\{1,2\},\{3,4\},\dots,\{2n-1,2n\}$ and apply pigeonhole.
\end{example}

\section{Inclusion--exclusion}

\begin{theorem}\label{thm:incexc}
For finite sets $A_1, \ldots, A_n$:
\[
\Bigl|\bigcup_{i} A_i\Bigr|
= \sum_i |A_i| - \sum_{i<j} |A_i \cap A_j|
+ \sum_{i<j<k} |A_i \cap A_j \cap A_k| - \cdots
+ (-1)^{n+1} |A_1 \cap \cdots \cap A_n|.
\]
\end{theorem}

\begin{example}[Derangements]\label{ex:derange}
A \emph{derangement} is a permutation with no fixed point. Let $A_i$ be the
permutations fixing $i$; inclusion--exclusion gives
\[
D_n = n!\sum_{k=0}^{n} \frac{(-1)^{k}}{k!}
\quad\xrightarrow[n \to \infty]{}\quad \frac{n!}{e}.
\]
So if $n$ guests grab random hats, the chance nobody gets their own tends to
$1/e \approx 0.3679$.
\end{example}

\section{Two classical identities}

\begin{theorem}[Vandermonde]\label{thm:vandermonde}
$\displaystyle\binom{m+n}{r} = \sum_{k=0}^{r} \binom{m}{k}\binom{n}{r-k}$
--- choose $r$ people from $m$ women and $n$ men by cases on how many women.
\end{theorem}

\begin{theorem}[Hockey stick]\label{thm:hockey}
$\displaystyle\sum_{i=r}^{n} \binom{i}{r} = \binom{n+1}{r+1}$
--- classify the $(r+1)$-subsets of $\{0,\dots,n\}$ by their largest
element.
\end{theorem}

\begin{example}[Poker hands]
From a $52$-card deck, $\binom{52}{5} = 2598960$ five-card hands. Full
houses: choose the triple rank, its suits, the pair rank, its suits:
$13 \cdot \binom{4}{3} \cdot 12 \cdot \binom{4}{2} = 3744$.
\end{example}

\section{Strategy: choosing the right formula}

Ask in order: (1) Am I \emph{selecting} or \emph{arranging} (does order
matter)? (2) Is repetition allowed? --- this places you in the four-problem
table. (3) Are objects distinguishable? If types repeat, use multiset
permutations or stars and bars. (4) ``At least/at most'' constraints often
yield to complement counting or inclusion--exclusion; existence claims to
pigeonhole.

\section{Summary}

Sum and product rules compose all elementary counts. Order + no repetition:
$P(n,k)$; no order: $\binom{n}{k}$; with repetition: $n^k$ and
$\binom{n+k-1}{k}$. Pascal's identity builds the triangle; the binomial and
multinomial theorems expand powers; pigeonhole guarantees collisions;
inclusion--exclusion corrects overcounting --- derangements being its
signature application.

\section*{Exercises}
\addcontentsline{toc}{section}{Exercises}

\begin{exercise}{\easy}
License plates consist of $2$ letters followed by $3$ digits. How many are
possible? How many with no repeated symbol?
\end{exercise}

\begin{exercise}{\easy}
Compute $P(7,3)$, $\binom{7}{3}$, and the number of binary strings of length
$7$ containing exactly three $1$s.
\end{exercise}

\begin{exercise}{\easy}
How many distinct arrangements has the word ``ABRACADABRA''?
\end{exercise}

\begin{exercise}{\medium}
A committee of $6$ is chosen from $10$ men and $8$ women. How many
committees contain (a) exactly $3$ women; (b) at least one woman?
\end{exercise}

\begin{exercise}{\medium}
How many solutions in nonnegative integers has $x + y + z = 11$? How many
with $x \geq 2$?
\end{exercise}

\begin{exercise}{\medium}
Find the coefficient of $x^{5}$ in $(2x - 3)^{8}$.
\end{exercise}

\begin{exercise}{\medium}
In how many ways can $8$ people sit around a round table? What if two
specific people insist on sitting together?
\end{exercise}

\begin{exercise}{\medium}
Show that among any $n + 1$ integers from $\{1, 2, \ldots, 2n\}$, one
divides another. \emph{Hint: write each as $2^a m$ with $m$ odd, and apply
pigeonhole to the odd parts.}
\end{exercise}

\begin{exercise}{\medium}
How many integers in $\{1, \ldots, 1000\}$ are divisible by $2$, $3$ or
$5$?
\end{exercise}

\begin{exercise}{\hard}
Prove Vandermonde's identity (Theorem~\ref{thm:vandermonde}) by computing
the coefficient of $x^{r}$ in $(1+x)^{m}(1+x)^{n} = (1+x)^{m+n}$.
\end{exercise}

\begin{exercise}{\hard}
Four guests check four hats and receive them back at random. Using
Example~\ref{ex:derange}, count the outcomes where nobody receives their own
hat, and the probability of that event.
\end{exercise}

\begin{exercise}{\hard}
How many surjections are there from a $5$-element set onto a $3$-element
set? \emph{Hint: inclusion--exclusion over the missed elements:}
$\sum_{j} (-1)^{j}\binom{3}{j}(3-j)^{5}$.
\end{exercise}

\section*{Solutions to the Exercises}
\addcontentsline{toc}{section}{Solutions to the Exercises}

\begin{solution}{10.1}
$26^2 \cdot 10^3 = 676000$. Without repetition:
$26 \cdot 25 \cdot 10 \cdot 9 \cdot 8 = 468000$.
\end{solution}

\begin{solution}{10.2}
$P(7,3) = 7\cdot6\cdot5 = 210$; $\binom{7}{3} = 35$; a string is determined
by the positions of its three $1$s: $\binom{7}{3} = 35$.
\end{solution}

\begin{solution}{10.3}
Letters: A$\times$5, B$\times$2, R$\times$2, C, D --- total $11$:
$\frac{11!}{5!\,2!\,2!\,1!\,1!} = 83160$.
\end{solution}

\begin{solution}{10.4}
(a) $\binom{8}{3}\binom{10}{3} = 56 \cdot 120 = 6720$.
(b) Total minus all-male: $\binom{18}{6} - \binom{10}{6} = 18564 - 210 =
18354$.
\end{solution}

\begin{solution}{10.5}
Stars and bars ($n = 3$, $k = 11$): $\binom{13}{2} = 78$. With $x \geq 2$,
substitute $x' = x - 2$: $x' + y + z = 9$, giving $\binom{11}{2} = 55$.
\end{solution}

\begin{solution}{10.6}
General term $\binom{8}{k}(2x)^{8-k}(-3)^{k}$; $x^5$ needs $k = 3$:
$\binom{8}{3} 2^{5} (-3)^{3} = 56 \cdot 32 \cdot (-27) = -48384$.
\end{solution}

\begin{solution}{10.7}
$(8-1)! = 5040$. Glue the pair into one block: $(7-1)! = 720$ circular
orders, times $2$ internal arrangements: $1440$.
\end{solution}

\begin{solution}{10.8}
Write each chosen integer as $2^{a} m$ with odd part
$m \in \{1, 3, \ldots, 2n-1\}$ --- only $n$ possible odd parts for $n+1$
numbers, so two share the same $m$: say $2^{a}m$ and $2^{b}m$ with $a < b$.
Then $2^{a}m \mid 2^{b}m$.
\end{solution}

\begin{solution}{10.9}
With $D_k$ = multiples of $k$: $|D_2| = 500$, $|D_3| = 333$, $|D_5| = 200$,
$|D_6| = 166$, $|D_{10}| = 100$, $|D_{15}| = 66$, $|D_{30}| = 33$.
Inclusion--exclusion: $500 + 333 + 200 - 166 - 100 - 66 + 33 = 734$.
\end{solution}

\begin{solution}{10.10}
The coefficient of $x^{r}$ in $(1+x)^{m+n}$ is $\binom{m+n}{r}$. In the
product $(1+x)^m (1+x)^n = \bigl(\sum_k \binom{m}{k}x^k\bigr)
\bigl(\sum_j \binom{n}{j}x^j\bigr)$, the $x^{r}$ terms pair $k + j = r$:
coefficient $\sum_{k=0}^{r} \binom{m}{k}\binom{n}{r-k}$. Equate.
\end{solution}

\begin{solution}{10.11}
$D_4 = 4!\bigl(1 - 1 + \tfrac12 - \tfrac16 + \tfrac1{24}\bigr) = 24 \cdot
\tfrac{9}{24} = 9$ derangements out of $4! = 24$ outcomes: probability
$\frac{9}{24} = \frac38 = 0.375$ (already close to $1/e$).
\end{solution}

\begin{solution}{10.12}
$\sum_{j=0}^{3} (-1)^{j} \binom{3}{j} (3-j)^{5}
= 3^5 - 3 \cdot 2^5 + 3 \cdot 1^5 - 0 = 243 - 96 + 3 = 150$.
\end{solution}
